% !PRINTED 676
\weberpage{676}
%<<CONT
und erhalten so aus (12), wenn $\lambda_s', \lambda_s'', \dots$ die conjugirten Logarithmen von $\eta', \eta'', \dots$ sind:
\weq{14}{
\begin{aligned}
a &< \Sigma g_s \lambda_s < a'\\
a' &< \Sigma g_s \lambda_s' < a''\\
a'' &< \Sigma g_s \lambda_s'' < a'''\\
&\cdots\cdots\cdots\cdots,
\end{aligned}
}
und die Reihe dieser Ungleichungen l\"asst sich unbegrenzt fortsetzen.

Alle diese Zahlen $\eta, \eta' \eta'', \dots$ haben \"uberdies nach (6) und (7) die Eigenschaft, dass ihre absoluten Normen $N, N', N'', \dots$ unter einer bestimmten endlichen Grenze $e^{nk}$ bleiben.

% [sic] The source omits the comma between \eta' and \eta'' in the preceding list.

Die Anzahl der m\"oglichen Werthe von $N, N', N'', \dots$ ist also endlich, n\"amlich gewiss nicht gr\"osser, als die gr\"osste in $e^{nk}$ enthaltene ganze Zahl. Ebenso ist die Anzahl aller m\"oglichen Reste der Zahlen $x_1, x_2, \dots, x_n$ [in (4)] nach allen Moduln $N, N', N'', \dots$ nur eine endliche, und daraus folgt, dass, wenn wir die Reihe der Zahlen (13) nur weit genug fortsetzen, in der Reihe \gesperrt{zwei verschiedene Zahlen $\eta^{(h)}, \eta^{(k)}$ auftreten m\"ussen, in denen $N^{(h)} = N^{(k)}$ und die Reste von $x_1, x_2, \dots, x_n$ nach dem Modul $N^{(h)}$ genau dieselben sind.}

Aus (14) aber folgt, wenn $h > k$ vorausgesetzt wird:
\weq{15}{\Sigma g_s (\lambda_s^{(h)} - \lambda_s^{(k)}) > 0.}
Ist $N$ die absolute Norm dieser beiden Zahlen $\eta^{(h)}, \eta^{(k)}$, so ist wegen der Uebereinstimmung der Reste der $x$ die Differenz $\eta^{(h)} - \eta^{(k)}$ durch $N$ theilbar. Andererseits ist nach \S.\,138\,(13) die Zahl $N$ durch $\eta^{(h)}$ theilbar, und folglich ist auch $\eta^{(h)} - \eta^{(k)}$ durch $\eta^{(h)}$ theilbar. Daraus ergiebt sich, dass auch $\eta^{(k)}$ durch $\eta^{(h)}$ und ebenso $\eta^{(h)}$ durch $\eta^{(k)}$ theilbar ist. Beide Zahlen sind also associirt, und wenn wir
\[
\frac{\eta^{(h)}}{\eta^{(k)}} = \varepsilon
\]
setzen, so ist $\varepsilon$ eine \gesperrt{Einheit} des K\"orpers $\Omega$.

Setzen wir ausserdem, indem wir mit $|\varepsilon_s|$ den absoluten Werth von $\varepsilon_s$ bezeichnen,
\weq{16}{\delta_s \log|\varepsilon_s| = l_s,}
so ergiebt sich aus (15):
\weq{17}{g_1 l_1 + g_2 l_2 + \cdots + g_\nu l_\nu > 0.}
